\chapter{Capítulo 2}

\section{Concepto de Ánalisis del puesto de trabajo}
Los puestos de trabajo están formados por una serie de tareas, obligaciones 
y responsabilidades. Las tareas son los elementos más básicos del trabajo; las obligaciones
 están formadas por una o más tareas que constituyen una actividad significativa 
 y la responsabilidad está definida por una o varias obligaciones que describen la razón de ser del trabajo.

 Para ello se realiza el llamado "análisis de puestos de trabajo", que se define
 como un proceso sistemático de recopilación de información para tomar decisiones relativas al trabajo,
 identificando las tareas, obligaciones y responsabilidad del puesto.

\section{Proceso de Realización del Ánalisis del puesto de trabajo}
Compuesto por 4 etapas:
\begin{enumerate}[label=(\Roman*)]
    \item Objetivos del análisis
    \item Delimitación del análisis
    \item Análisis de puestos
    \item Descripción y especificación del puesto
\end{enumerate}

\subsection{Etapa 1: Objetivos del análisis}
Es necesario identificar la información que se desea obtener y el procedimiento a seguir. Para ello
se recopila la información básica relevante sobre el puesto.
\subsection{Etapa 2: Delimitación del análisis}
Se selecciona que puestos a analizar, cuándo y quién se encargará de ello.
¿Qué puestos se van a analizar?¿Cuándo se hará el estudio?¿Quién lo realizará?

\subsection{Etapa 3: Análisis de puestos}
Reunir datos sobre conductas, condiciones y capacidades requeridad para el puesto.
\begin{itemize}
    \item \underline{Entrevistas}: Permite obtener información directa, pueden ser individuales, de grupo. Es costosa y requiere tiempo.
    \item \underline{Observación}: Observar al trabajador trabajando. Es lenta y costosa.
    \item \underline{Diarios o bitácotras}: Registro llevado por el trabajador.
    \item \underline{Cuestionarios}: El trabajador responde una serie de preguntas sobre el puesto. Requiere tiempo de preparación y comprobación.
    \item \underline{Grupos de expertos}: Se recoge información de un grupo de especialistas, normalmente trabajadores. Lento y costoso.
\end{itemize}

Técnicas de análisis:
\begin{itemize}
    \item \underline{Análisis del inventario de tareas}
\end{itemize}
